% Pursuing Stacks (cotes 134-2 to 134-8): archive cote 134-4, site page 58 (PDF page 59), typescript page 146.
% typescript. Draft compared: ps3.tex:2956-3022.
% AI-transcribed by Claude Opus 5.5 (claude-opus-5-5, adaptive thinking, effort high), Message Batches API, batch ps2 (msgbatch_01NVwr1Yuo8qJf5rykeGBXPK), 2026-11-12; not yet checked by eye.
\scanpage{134-4/58}
\mspage{146}

From our point of view, the main point for paying attention to the
homotopy relation for functors, is of course because homotopic functors
define the same map in the localized category (Hot), which\inserted{category} for the time
being is (together with the modelizers designed for describing it) our
main focus of attention. According to the general scheme of homotopy
theory as reviewed previously, homotopy classes of functors give rise
to a quotient category of (Cat), which is at the same time a localization
of (Cat) with respect to homotopisms, and which we'll denote by $\overline{\text{(Cat)}}$.
Thus we get a factorisation of the canonical functor fro\hand{m}(Cat) to (Hot)
\[ \text{(Cat)} \hand{$\longrightarrow$} \overline{\text{(Cat)}} \hand{$\longrightarrow$} \text{(Hot)} \quad . \]
Th\uline{is} is just the factorisation of the localization functor with respect
to $W_{\text{(Cat)}} = W$ , through the ``\uline{partial}\xout{ly} \uline{localization}'' with respect to the
smaller set $W_h \hand{$\subset$} W$ consisting of homotopi\unsure{sms}\hand{|}only. All we have used about
$W = W_{\text{(Cat)}}$ (from which the localization (Hot) $= W^{-1}$(Cat) is deduced) for
getting this factorization, was that $f \hand{$\in$} W$ implies $\hand{$\pi$}_0(f)$ bijective,\hand{*} and
\margin{*6.4. As a matter of fact, the condition about $\pi_0$ visibly \unsure{is'nt|isn't} needed!}
that the projections $C\mathrm{x}\hand{$\Delta$}_1 \hand{$\longrightarrow$} C$ are in W. (Of course, we assume tacitly
that W is saturated too.) \xout{Well known analogies would} We may view (Hot)
as a localization of $\overline{\text{(Cat)}}$ with respect to the set of arrows $\overline{W}$ correspon%
\xout{Well known analogies would suggest at this point}
ding to W - as a matter of fact, W may be viewed as the inverse image
of $\overline{W}$ by the canonical functor from (Cat) to it's partial localization
$\overline{\text{(Cat)}}$ - the category of ``homotopy types'' relative to the canonical
homotopy structure of (Cat).

Well known analogies would suggest at this point that we may well be
able to describe (Hot) in terms of ( $\overline{\text{(Cat)}}$, $\overline{W}$ ) by a \uline{calculus of frac-}%
\uline{tions} - right fractions presumably, or maybe either right or left. This
may possibly lead to a direct proof of the notion of weak equivalence we
have been working with being strongly saturated, without having to rely
upon Quillen's closed model theory. But it is not yet the moment to pursue
this line of thought, which would take us off the main focus\xout{x} at present.

59. Before pushing ahead, I would like to make still another point
about the work done yesterday - a point suggested by looking at the
case of the modelizer (Spaces) , which after all is the next best
``naive'' modelizer, less close to algebra than (Cat), but still worth
being taken into account ! This category satisfies the conditions a) to
f) , \uline{except} the condition c) - which would mean that the connected com-%
ponents of a space (as defined in terms of ususal topology) are open sub-%
sets, which is true (for a space and it's open subsets) \uline{only} for locally
connected spaces. The point is that this does'nt (or should'nt) really

\margin{58}
